% GENERATED FILE. Edit canonical lessons, then run npm run generate:books.
% canonical-source-hash: fnv1a64:70630939d6b296b8
% canonical-lessons: LA-C401-paedagogus, LA-C401-praeceptor, LA-C401-grammaticus, LA-C401-litterator, LA-C401-condiscipulus

\chapter{At School}
\label{ch:la-at-school}
\hlchaptermodality{401}

\begin{chapteropening}
\textbf{By the end of this chapter:} \emph{I can name the slave who walks a boy to school, a teacher, the grammar master, the primary teacher and a schoolfellow.}

Complete the last lesson of chapter 401: I can name the slave who walks a boy to school, a teacher, the grammar master, the primary teacher and a schoolfellow.
\end{chapteropening}

\section[paedagōgus, paedagōgī]{paedagōgus, paedagōgī — a slave who takes a boy to school; a tutor}

\label{lesson:LA-C401-paedagogus}



\begin{quote}
Before the new one: say the Latin for to equip, then the Latin for to instruct.
\end{quote}

\subsection*{You'll want to know: paedagōgus, paedagōgī}
\textbf{paedagōgus, paedagōgī} — \textquotedblleft{}a slave who takes a boy to school; a tutor\textquotedblright{}.

\begin{quote}
Paedagōgus puerum ad scholam dūcit.
\end{quote}

\begin{quote}
\emph{The slave-tutor takes the boy to school.}
\end{quote}

\emph{paedagōgus, paedagōgī} m., borrowed from Greek \emph{paidagōgos}, child-leader. English \emph{pedagogue}.

\begin{cousinweb}
\emph{paedagōgus} was borrowed from Greek \emph{paidagōgos} \textquotedblleft{}child-leader\textquotedblright{} (EtymDB), built from \emph{pais} \textquotedblleft{}a child\textquotedblright{} and \emph{agōgos} \textquotedblleft{}leading\textquotedblright{}. English \emph{pedagogue} comes from it.
\end{cousinweb}

\begin{grammarlens}[title={What you've built}]
One more word for telling, explaining and describing at B1.
\end{grammarlens}

\subsection*{Your turn}
\begin{itemize}
\raggedright
  \item \emph{Say it:} \emph{paedagōgus, paedagōgī}
  \item \emph{Say it:} \emph{paedagōgus, paedagōgī}, once more
  \item \emph{Say it:} a sentence of your own with \emph{paedagōgus, paedagōgī}
  \item \emph{Recall:} say the Latin for to equip, then the Latin for to instruct, then say \emph{paedagōgus, paedagōgī} again
\end{itemize}

\begin{tcolorbox}[breakable,colback=teal!4,colframe=teal!35!black,title={Before you move on}]
What does paedagōgus, paedagōgī mean? (\textquotedblleft{}A slave who takes a boy to school; a tutor\textquotedblright{}.) Say it once more.
\end{tcolorbox}
\section[praeceptor, praeceptōris]{praeceptor, praeceptōris — a teacher, an instructor}

\label{lesson:LA-C401-praeceptor}



\begin{quote}
Before the new one: say the Latin for to instruct, then the Latin for a slave who takes a boy to school.
\end{quote}

\subsection*{You'll want to know: praeceptor, praeceptōris}
\textbf{praeceptor, praeceptōris} — \textquotedblleft{}a teacher, an instructor\textquotedblright{}.

\begin{quote}
Praeceptor discipulōs suōs laudāvit.
\end{quote}

\begin{quote}
\emph{The teacher praised his pupils.}
\end{quote}

\emph{praeceptor, praeceptōris} m., built on \emph{praecipiō}: the one who lays down rules. English \emph{preceptor}.

\begin{cousinweb}
\emph{praeceptor} is built on \emph{praecipiō} \textquotedblleft{}to instruct\textquotedblright{} (Lewis \& Short), with the ending \emph{-tor} of the one who does it.
\end{cousinweb}

\begin{grammarlens}[title={What you've built}]
One more word for telling, explaining and describing at B1.
\end{grammarlens}

\subsection*{Your turn}
\begin{itemize}
\raggedright
  \item \emph{Say it:} \emph{praeceptor, praeceptōris}
  \item \emph{Say it:} \emph{praeceptor, praeceptōris}, once more
  \item \emph{Say it:} a sentence of your own with \emph{praeceptor, praeceptōris}
  \item \emph{Recall:} say the Latin for to instruct, then the Latin for a slave who takes a boy to school, then say \emph{praeceptor, praeceptōris} again
\end{itemize}

\begin{tcolorbox}[breakable,colback=teal!4,colframe=teal!35!black,title={Before you move on}]
What does praeceptor, praeceptōris mean? (\textquotedblleft{}A teacher, an instructor\textquotedblright{}.) Say it once more.
\end{tcolorbox}
\section[grammaticus, grammaticī]{grammaticus, grammaticī — a teacher of literature and grammar}

\label{lesson:LA-C401-grammaticus}



\begin{quote}
Before the new one: say the Latin for a slave who takes a boy to school, then the Latin for a teacher.
\end{quote}

\subsection*{You'll want to know: grammaticus, grammaticī}
\textbf{grammaticus, grammaticī} — \textquotedblleft{}a teacher of literature and grammar\textquotedblright{}.

\begin{quote}
Apud grammaticum poētās legēbāmus.
\end{quote}

\begin{quote}
\emph{With the grammar master we read the poets.}
\end{quote}

\emph{grammaticus, grammaticī} m., borrowed from Greek \emph{grammatikos}. After the first school, a boy went to the \emph{grammaticus} to read the poets.

\begin{cousinweb}
\emph{grammaticus} was borrowed from Greek \emph{grammatikos} \textquotedblleft{}skilled in letters\textquotedblright{} (EtymDB), built on \emph{gramma} \textquotedblleft{}a letter\textquotedblright{}. English \emph{grammar} comes from the same Greek family.
\end{cousinweb}

\begin{grammarlens}[title={What you've built}]
One more word for telling, explaining and describing at B1.
\end{grammarlens}

\subsection*{Your turn}
\begin{itemize}
\raggedright
  \item \emph{Say it:} \emph{grammaticus, grammaticī}
  \item \emph{Say it:} \emph{grammaticus, grammaticī}, once more
  \item \emph{Say it:} a sentence of your own with \emph{grammaticus, grammaticī}
  \item \emph{Recall:} say the Latin for a slave who takes a boy to school, then the Latin for a teacher, then say \emph{grammaticus, grammaticī} again
\end{itemize}

\begin{tcolorbox}[breakable,colback=teal!4,colframe=teal!35!black,title={Before you move on}]
What does grammaticus, grammaticī mean? (\textquotedblleft{}A teacher of literature and grammar\textquotedblright{}.) Say it once more.
\end{tcolorbox}
\section[litterātor, litterātōris]{litterātor, litterātōris — an elementary teacher of letters}

\label{lesson:LA-C401-litterator}



\begin{quote}
Before the new one: say the Latin for a teacher, then the Latin for a teacher of literature and grammar.
\end{quote}

\subsection*{You'll want to know: litterātor, litterātōris}
\textbf{litterātor, litterātōris} — \textquotedblleft{}an elementary teacher of letters\textquotedblright{}.

\begin{quote}
Litterātor puerōs legere et scrībere docet.
\end{quote}

\begin{quote}
\emph{The elementary teacher teaches the boys to read and write.}
\end{quote}

\emph{litterātor, litterātōris} m., built on \emph{littera}, a letter. He was also called \emph{lūdī magister}.

\begin{cousinweb}
\emph{litterātor} is built on \emph{littera} \textquotedblleft{}a letter\textquotedblright{} (EtymDB; Lewis \& Short), with the ending \emph{-tor} of the one who does it.
\end{cousinweb}

\begin{grammarlens}[title={What you've built}]
One more word for telling, explaining and describing at B1.
\end{grammarlens}

\subsection*{Your turn}
\begin{itemize}
\raggedright
  \item \emph{Say it:} \emph{litterātor, litterātōris}
  \item \emph{Say it:} \emph{litterātor, litterātōris}, once more
  \item \emph{Say it:} a sentence of your own with \emph{litterātor, litterātōris}
  \item \emph{Recall:} say the Latin for a teacher, then the Latin for a teacher of literature and grammar, then say \emph{litterātor, litterātōris} again
\end{itemize}

\begin{tcolorbox}[breakable,colback=teal!4,colframe=teal!35!black,title={Before you move on}]
What does litterātor, litterātōris mean? (\textquotedblleft{}An elementary teacher of letters\textquotedblright{}.) Say it once more.
\end{tcolorbox}
\section[condiscipulus, condiscipulī]{condiscipulus, condiscipulī — a schoolfellow, a fellow pupil}

\label{lesson:LA-C401-condiscipulus}



\begin{quote}
Before the new one: say the Latin for a teacher of literature and grammar, then the Latin for an elementary teacher of letters.
\end{quote}

\subsection*{You'll want to know: condiscipulus, condiscipulī}
\textbf{condiscipulus, condiscipulī} — \textquotedblleft{}a schoolfellow, a fellow pupil\textquotedblright{}.

\begin{quote}
Condiscipulī in eādem scholā sedēbant.
\end{quote}

\begin{quote}
\emph{The schoolfellows sat in the same school.}
\end{quote}

\emph{condiscipulus, condiscipulī} m., \emph{con-} with + \emph{discipulus}.

\begin{cousinweb}
\emph{condiscipulus} is built from \emph{con-} \textquotedblleft{}with\textquotedblright{} and \emph{discipulus} \textquotedblleft{}a pupil\textquotedblright{} (Lewis \& Short).
\end{cousinweb}

\begin{grammarlens}[title={What you've built}]
One more word for telling, explaining and describing at B1.
\end{grammarlens}

\subsection*{Your turn}
\begin{itemize}
\raggedright
  \item \emph{Say it:} \emph{condiscipulus, condiscipulī}
  \item \emph{Say it:} \emph{condiscipulus, condiscipulī}, once more
  \item \emph{Say it:} a sentence of your own with \emph{condiscipulus, condiscipulī}
  \item \emph{Recall:} say the Latin for a teacher of literature and grammar, then the Latin for an elementary teacher of letters, then say \emph{condiscipulus, condiscipulī} again
\end{itemize}

\begin{tcolorbox}[breakable,colback=teal!4,colframe=teal!35!black,title={Before you move on}]
What does condiscipulus, condiscipulī mean? (\textquotedblleft{}A schoolfellow, a fellow pupil\textquotedblright{}.) Say it once more.
\end{tcolorbox}
